% ============================================================================
\section{Conclusion}
\label{sec:conclusion}
% ============================================================================

The transverse rate underlying microstructure imaging is an intrinsic bulk rate plus a
surface-relaxivity rate $\rho\,(S/V)$ --- both isotropic and TE-linear, hence
non-identifiable, with $S/V$ unresolved. It is the relaxometric face of the wall collisions
diffusion MRI reads as time-dependent extra-axonal diffusion, and it acts with no gradient;
a wall-counting Monte Carlo reproduces it from first principles. Because the interior and
exterior walls carry different $S/V$, the compartments relax differently, and any estimate
that normalises by a TE-weighted $b{=}0$ is biased. On the diffusion side this over-weights
the intra-axonal signal (the leading-order intra-axonal fraction bias) on physiologically
dense white matter --- by ${\approx}12\%$ over the robust packing band at clinical echo
time and the cited $\rho$, with a sign set by a single packing ratio
$(1-\mathrm{VF})/(g\,\mathrm{VF})$ that crosses unity at $\mathrm{VF}^{\ast}=1/(1+g)$, and a
testable, packing-dependent TE drift --- a first-order, systematic bias on a headline metric
that shifts regional and cohort means rather than averaging out. The same physics reads through relaxometry as a smaller, structural
myelin-water bias: the thinnest axons' water is dragged below the myelin window and counted
as myelin, so fine white matter reads myelin-\emph{richer} at equal myelin content
(${\sim}0.33$\,pp, ${\sim}0.82$\,pp of inferred myelin after the $1/W_M$ amplification,
climbing super-linearly at the upper end of the order-of-magnitude-uncertain $\rho$); being
lumen-set, it \emph{cancels} in the within-subject longitudinal change that primary
demyelination traces. Diffusion is in principle the orthogonal axis that could constrain
the geometry directly --- the same wall collisions, read with a gradient --- but realising
it is gradient-limited, left to future work.
